% !TEX root = ../main.tex
\documentclass[../main.tex]{subfiles}

\begin{document}

\selectlanguage{english}

\begin{center}
    \textbf{\LARGE EXECUTIVE ARTIFICIAL INTELLIGENCE ESSAY} \\
    \vspace{0.2cm}
    \textbf{\large Topic: Multiple Drone Attack in Urban City} \\
    \textit{(Strategic Defense Optimization Against Multi-Target UAV Attacks in Urban Environments)}
\end{center}

\vspace{0.5cm}

\section{PART 1: PROBLEM STATEMENT \& OBJECTIVE (15\% Weight)}

\begin{itemize}
    \item \textbf{Real-World Context:} In modern security paradigms, the proliferation of low-cost Unmanned Aerial Vehicles (drones) presents an asymmetric threat to urban security \cite{mutzari2025defending}. Malicious actors can deploy coordinated multi-drone attacks targeting critical infrastructure such as government buildings, power grids, water supply networks, and hospitals \cite{mutzari2025defending}. Urban topographies, characterized by high-rise buildings, multi-layered obstacles, and signal interference, completely neutralize traditional ground-based radar and defensive hardware \cite{mutzari2025defending}. Consequently, deploying a highly mobile, aerial interceptor squadron of defense drones is imperative for urban protection \cite{mutzari2025defending}.
    \item \textbf{Core Objective:} This project rejects the passive, reactive ``catch-as-catch-can'' approach of legacy tracking systems \cite{mutzari2025defending}. Instead, the objective is to implement strategic AI frameworks designed to \underline{minimize total systemic damage} to human lives and physical assets by optimizing defense resource allocation \cite{mutzari2025defending}. The framework models the problem under the operational constraint that defensive assets are inherently limited compared to the adversary's offensive capabilities \cite{mutzari2025defending}.
\end{itemize}

\end{document}
